\documentclass[11pt, a4paper]{amsart}
\usepackage{amsmath, amssymb, amsthm, amsfonts}
\usepackage{mathtools}
\usepackage[margin=1in]{geometry}
\usepackage{hyperref}

\newtheorem{theorem}{Theorem}[section]
\newtheorem{lemma}[theorem]{Lemma}
\newtheorem{proposition}[theorem]{Proposition}
\newtheorem{corollary}[theorem]{Corollary}
\newtheorem{fact}[theorem]{Standard Fact}
\theoremstyle{definition}
\newtheorem{definition}[theorem]{Definition}
\newtheorem{remark}[theorem]{Remark}
\newtheorem{notation}[theorem]{Notation}
\newtheorem{hypothesis}[theorem]{Hypothesis}

\newcommand{\R}{\mathbb{R}}
\newcommand{\Q}{\mathbb{Q}}
\newcommand{\Z}{\mathbb{Z}}
\newcommand{\N}{\mathbb{N}}
\newcommand{\id}{\operatorname{id}}
\newcommand{\tr}{\operatorname{tr}}
\newcommand{\spec}{\operatorname{spec}}
\newcommand{\op}{\operatorname{op}}
\newcommand{\Sym}{\operatorname{Sym}}
\newcommand{\diam}{\operatorname{diam}}
\newcommand{\vol}{\operatorname{vol}}
\newcommand{\dGH}{d_{\mathrm{GH}}}
\newcommand{\dsph}{d_{S^3}}
\newcommand{\supp}{\operatorname{supp}}

\title{Cascade--Classical Correspondence VI:\\
  Schl\"{a}fli Gromov--Hausdorff Convergence\\
  (with a conditional spectral-convergence programme)}
\author{}
\date{}

\begin{document}
\maketitle

\begin{abstract}
We prove quantitative Gromov--Hausdorff convergence of the
$D_4$-branch Schl\"{a}fli refinement family
\cite{RefinementSpaces} to the round unit $3$-sphere
$(S^3, g_{\rm round})$. We equip $V_n^{D_4}$ with an
\emph{intrinsic graph metric} $d_n$ whose edge lengths are
explicit spherical arc lengths under an iterated-spherical-midpoint
embedding $\iota_n \colon V_n^{D_4} \to S^3$, and prove:
(T1) \emph{conditional on Hypothesis~\ref{fact:tiling}} (spherical
cells tile $S^3$), $\varepsilon_n = O(h_n)$ net radius, and
$\iota_n$ is an isometric embedding of the extrinsic metric
$(V_n^{D_4}, d_n^{\rm ext})$ into $(S^3, \dsph)$;
(T2) \emph{conditional on Hypothesis~\ref{fact:tiling}},
$\dGH((V_n^{D_4}, d_n^{\rm ext}), (S^3, \dsph)) = O(h_n)$, immediate
from T1; the mesh $h_n$ satisfies $h_n \leq M_\ast \alpha^n h_0$
for any $\alpha \in (1/\sqrt{3}, 1)$ and an absolute constant
$M_\ast = M_\ast(\alpha) \geq 1$
(Lemma~\ref{lem:induction-closes}), so the GH rate is $O(\alpha^n)$
for any such $\alpha$ up to an absolute multiplicative constant
(not the dyadic $2^{-n}$ rate with constant $1$);
(Prop.~\ref{prop:doubling}) doubling with explicit constants
from independent covering and separation bounds,
\emph{conditional additionally on Hypothesis~\ref{fact:separation}}
(uniform edge-length separation).

On the spectral side, we prove the $V_0$-vertex-figure isotropy
of the $D_4$ Schl\"{a}fli refinement: the weighted outer-product
sum $\sum (\ell_e^{(n)}/h_n)^2\,\widehat\tau_e \otimes
\widehat\tau_e$ at every level-$0$ orbit representative equals
$(8/3)\,I_{T_v S^3}$ exactly
(Proposition~\ref{prop:v0-at-level-n}), identifying $8/3$ as the
structural constant governing graph-Laplacian-to-Laplace-Beltrami
convergence.

We also present, as a \emph{conditional program}, a
Courant--Fischer-based spectral-convergence argument under
Hypothesis~\ref{hyp:isotropy} (uniform $8/3$ weighted-isotropy
with $O(h_n^2)$ error at every vertex, verified unconditionally
for $V_0$; extension to midpoint/centroid orbits deferred).
Theorems~\ref{thm:form-consistency}
and~\ref{thm:spectral-convergence} state this program; proofs
of several load-bearing lemmas (quadrature
Lemma~\ref{lem:quadrature}, dual form consistency
Lemma~\ref{lem:dual-form-consistency}) are sketched, and the
spectral statement is not yet Clay-grade theorem-level output.
This is the target of a follow-up paper.
We make no use of paper P4, no Lorentzian/Einstein statement, and
no continuum-PDE claim.
\end{abstract}

\tableofcontents

%=====================================================================
\section{Introduction and Source Audit}
\label{sec:intro}
%=====================================================================

\subsection{What this paper proves, and what is conditional}

\textbf{Theorems (each conditional on the indicated structural
hypothesis):}
\begin{description}
\item[\textbf{T1}] \emph{$\varepsilon$-net}
  (Theorem~\ref{thm:epsilon-net}, conditional on
  Hypothesis~\ref{fact:tiling}): $\varepsilon_n = O(h_n)$;
  $\iota_n$ is an isometric embedding of $(V_n, d_n^{\rm ext})$
  into $(S^3, \dsph)$ under the extrinsic metric.
  The intrinsic-vs-extrinsic bi-Lipschitz comparison
  (Theorem~\ref{thm:intrinsic-extrinsic}) is additionally
  \emph{conditional} on Hypothesis~\ref{hyp:isotropy}, as it
  requires a greedy pearl-string construction along geodesics
  that fails at non-$V_0$ orbits without uniform isotropy.
\item[\textbf{T2}] \emph{Quantitative GH convergence}
  (Theorem~\ref{thm:gh-convergence}, conditional on
  Hypothesis~\ref{fact:tiling}): $\dGH((V_n, d_n^{\rm ext}),
  (S^3, \dsph)) = O(h_n)$, with $h_n \leq M_\ast \alpha^n h_0$
  for any $\alpha \in (1/\sqrt{3}, 1)$ up to an absolute
  constant $M_\ast = M_\ast(\alpha)$; immediate from T1 via
  \cite[Prop.~7.3.10]{BuragoBuragoIvanov}.
\item[\textbf{Doubling}] (Proposition~\ref{prop:doubling},
  conditional on Hypotheses~\ref{fact:tiling} and~\ref{fact:separation}):
  explicit constants from separate covering (Theorem~\ref{thm:epsilon-net})
  and separation (Hypothesis~\ref{fact:separation}) bounds.
\item[\textbf{$V_0$-orbit isotropy}] (Proposition~\ref{prop:v0-at-level-n}):
  the vertex-figure tensor equals $(8/3)\,I_{T_v S^3}$ exactly,
  computed by direct enumeration, for every $v \in V_0$ at every
  level $n$. Cube geometry of the $24$-cell vertex figure gives
  the isotropy exactly, not approximately. (This statement does
  not require Hypotheses~\ref{fact:tiling} or~\ref{fact:separation};
  it is a level-$0$-orbit-only fact.)
\end{description}

\textbf{Conditional program (stated as a program toward spectral
convergence, not as proved theorems at Clay level):}
\begin{description}
\item[\textbf{T3 (program)}] \emph{Dirichlet-form consistency}
  (Theorem~\ref{thm:form-consistency}, proved at $C^2$ regularity
  under Hypothesis~\ref{hyp:isotropy} using the
  quadrature Lemma~\ref{lem:quadrature} whose full proof is
  sketched) and \emph{Eigenvalue
  convergence} (Theorem~\ref{thm:spectral-convergence}, proved
  via Courant--Fischer min-max using
  Lemma~\ref{lem:dual-form-consistency} whose full proof is
  also sketched):
  proved conditional on Hypothesis~\ref{hyp:isotropy}
  (approximate isotropy at every $V_n$ vertex with $O(h_n^2)$
  error). Hypothesis~\ref{hyp:isotropy} is PROVED for $V_0$
  orbit unconditionally; extending to midpoint and centroid
  orbits is the subject of a follow-up paper.
\end{description}

\textbf{Not claimed:}
Lorentzian limit; Einstein equations; continuum PDE; hydrodynamic
statements; $H_4$-branch analysis; direct dependence on paper P4;
BIK-style spectral convergence (their $\rho$-proximity framework
does not match the fixed Schl\"{a}fli graph; we prove the
spectral theorem directly).

\subsection{The structural insight: why the Schl\"{a}fli
  refinement converges to Laplace--Beltrami}
\label{subsec:story}

Not every $\varepsilon$-net of $S^3$ produces a graph Laplacian
that converges to $-\Delta_{S^3}$. Generically, the edge-direction
distribution at each vertex must be approximately \emph{isotropic}
for the discrete energy $\sum_{\{u, v\}} w_{uv}\,(f(u) - f(v))^2$
to approximate the isotropic Dirichlet energy
$\int |\nabla f|^2 d\mu$.

The $D_4$-branch Schl\"{a}fli refinement satisfies this isotropy
\emph{because of its symmetry group}. The $24$-cell has
automorphism group containing $W(D_4)$ acting transitively on
$V_{24}$, and every vertex of $V_{24}$ has a vertex figure equal
to a cube: the $8$ nearest neighbors form the $8$ vertices of a
cube in the tangent plane at that vertex. The outer-product sum
over a cube's vertices is exactly isotropic:
$\sum_{\varepsilon \in \{\pm 1\}^3}
(\varepsilon/\sqrt{3}) \otimes (\varepsilon/\sqrt{3}) = (8/3)\,
I_{\R^3}$. This identity, together with
Proposition~\ref{prop:v0-at-level-n}, identifies $8/3$ as the
structural constant of the $D_4$ Schl\"{a}fli refinement.

Under Hypothesis~\ref{hyp:isotropy} (the same constant, with
$O(h_n^2)$ error, at every $V_n$ vertex), the graph Laplacian
converges to Laplace--Beltrami at the correct scale. Without the
hypothesis, only the unconditional $V_0$-vertex isotropy
identity (Proposition~\ref{prop:v0-at-level-n}) holds; full
spectral convergence requires the hypothesis. (An earlier draft
mentioned a ``neighborhood of $V_0$'' fallback; since $V_0$ is
a finite discrete vertex set, not a neighborhood in $S^3$, that
fallback does not apply, and we withdraw it.)

The codex iter-2 blockers (``$d_n$ tautological'', ``T1(b)
arithmetic wrong'', ``$\Delta_n$ ignores edge lengths'', ``BIK
mismatch'') each pointed at an aspect of this story: that the
Schl\"{a}fli graph's spectral structure is determined by its
$W(D_4)$ equivariance, not by generic $\varepsilon$-net theory.
The present iteration encodes this directly.

\subsection{Source audit}

\textbf{Cascade-internal:}
P2~\cite{AlgebraicSubstrate} ($24$-cell coordinates, $W(D_4)$
action);
P3~\cite{RefinementSpaces} (refinement graphs, vertex inclusion);
P5~\cite{MetricProjection} (refined-vertex coordinate
equivariance). P4 NOT cited.

\textbf{Classical:}
\cite{BuragoBuragoIvanov} (GH and $\varepsilon$-isometry);
\cite{doCarmo, PetersenRiemannian} (spherical and Riemannian
geometry); \cite{BerardSpectralGeometry} ($S^3$ spectrum);
\cite[\S VIII.8]{Kato} (Courant--Fischer min-max);
\cite{ReedSimonI} (basic FA);
\cite{CoxeterRegularPolytopes} ($24$-cell and Schl\"{a}fli
refinement); \cite{Humphreys} (Coxeter groups);
\cite{HeinonenAnalysis} (doubling).
\cite{BuragoIvanovKurylev} cited for context only --- not used
as load-bearing machinery.

%=====================================================================
\section{The $D_4$ Branch and Base Geometry}
\label{sec:D4-base}
%=====================================================================

\begin{notation}[Level-$n$ vertices, coordinates, $S^3$ metric]
\label{not:basic}
$V_n := V(G_n^{D_4})$ for $n \geq 0$
\cite[Def.~2.1, Def.~2.2]{RefinementSpaces}. $x_v^{(n)} \in \R^4$
is the Euclidean coordinate from P5~\cite[Notation~2.1]{MetricProjection}.
$S^3 := \{x \in \R^4 : \|x\| = 1\}$, $\dsph(a, b) :=
\arccos(\langle a, b \rangle) \in [0, \pi]$.
\end{notation}

\begin{fact}[$V_0$ on unit sphere, $8$-regular, edges at $\pi/2$]
\label{fact:V0}
$V_0 = V_{24} = \{\pm e_i : 1 \leq i \leq 4\} \cup
\{\tfrac{1}{2}(\pm 1, \pm 1, \pm 1, \pm 1)\}$. All $v \in V_0$
have $\|x_v\| = 1$. The $1$-skeleton is $8$-regular. For every
$\{u, v\} \in E(G_0^{D_4})$, $\langle x_u, x_v\rangle = 0$, so
$\dsph(x_u, x_v) = \pi/2$.
\end{fact}

\begin{proof}
Explicit computation from the $24$-cell coordinates
\cite[\S 5]{AlgebraicSubstrate}. The $24$-cell has $24$ vertices,
edge length $\sqrt{2}$ on the unit sphere, inner product $0$ on
adjacent vertices, and each vertex has $8$ nearest neighbors
forming a cube in the tangent $3$-plane
\cite[\S 22.2]{CoxeterRegularPolytopes}.
\end{proof}

\begin{notation}[$W(D_4)$]
\label{not:WD4}
$W(D_4) \leq O(4)$ is the Weyl group of $D_4$
\cite[Ch.~1]{Humphreys}. $|W(D_4)| = 192$. $W(D_4)$ acts
transitively on $V_{24}$ (vertex stabilizer order $8$,
isomorphic to a hyperoctahedral subgroup). Coordinate
equivariance: $x_{g \cdot v}^{(n)} = g\,x_v^{(n)}$ for all
$g \in W(D_4)$, $v \in V_n$, by \cite[Lem.~2.4]{MetricProjection}.
\end{notation}

\begin{remark}[Why spherical midpoints, not Euclidean]
\label{rmk:why-spherical}
The Euclidean refinement of P3 produces refined vertices
that drift into the interior of the unit ball: the midpoint of
$\pi/2$-separated unit vectors has norm $1/\sqrt{2}$. Iterated
barycenters at arbitrary depth include points of arbitrarily
small norm, so naive radial projection
$\R^4 \setminus \{0\} \to S^3$ lacks uniform bi-Lipschitz control.
We use \emph{spherical} midpoints and centroids
(Definition~\ref{def:iota-n}) instead, keeping all refined
vertices on $S^3$.
\end{remark}

%=====================================================================
\section{The Spherical Realization}
\label{sec:spherical}
%=====================================================================

\begin{definition}[Spherical midpoint and centroid]
\label{def:sph-mid}
For $a_1, \ldots, a_k \in S^3$ with $S := a_1 + \cdots + a_k \neq
0$, the \emph{spherical centroid} is
$\mathrm{sph.cen}(a_1, \ldots, a_k) := S / \|S\|$. For $a, b \in
S^3$ with $a \neq -b$, the \emph{spherical midpoint} is
$\mathrm{sph.mid}(a, b) := (a + b)/\|a + b\|$.
\end{definition}

\begin{definition}[Spherical vertex map $\iota_n$]
\label{def:iota-n}
$\iota_n \colon V_n \to S^3$ inductively:
$\iota_0(v) := x_v^{(0)}$;
$\iota_{n+1}(v) := \iota_n(v)$ for $v \in V_n$;
$\iota_{n+1}(m_{\{u_1, u_2\}}) :=
\mathrm{sph.mid}(\iota_n(u_1), \iota_n(u_2))$ for edge
midpoints;
$\iota_{n+1}(c_F) :=
\mathrm{sph.cen}(\iota_n(u_1), \ldots, \iota_n(u_{k_F}))$ for
face centroids $F = \{u_1, \ldots, u_{k_F}\}$. Well-definedness
(non-antipodal inputs, nonvanishing Euclidean sums) is
verified in Lemma~\ref{lem:induction-closes}.
\end{definition}

\begin{definition}[Mesh scale $h_n$]
\label{def:hn}
$h_n := \max_{\{u, v\} \in E(G_n^{D_4})}
\dsph(\iota_n(u), \iota_n(v))$.
\end{definition}

\begin{lemma}[Spherical midpoint halves arc length]
\label{lem:sph-mid-shrink}
For $a, b \in S^3$ with $\dsph(a, b) = \theta \in [0, \pi)$,
$m := \mathrm{sph.mid}(a, b)$ has $\dsph(a, m) = \dsph(b, m) =
\theta/2$.
\end{lemma}

\begin{proof}
$\langle a + b, a\rangle = 1 + \cos\theta$, $\|a + b\|^2 = 2(1 +
\cos\theta) = 4\cos^2(\theta/2)$, so $\|a + b\| = 2\cos(\theta/2)$.
Then $\langle m, a\rangle = (1 + \cos\theta)/(2\cos(\theta/2)) =
\cos(\theta/2)$, giving $\dsph(a, m) = \theta/2$.
\end{proof}

\begin{lemma}[Inductive closure and mesh decay]
\label{lem:induction-closes}
Definition~\ref{def:iota-n} is well-defined at every level.
For every $\alpha \in (1/\sqrt{3}, 1)$ (for instance $\alpha =
0.6$), there exists $N_\ast = N_\ast(\alpha)$ such that
$h_{n+1} \leq \alpha\,h_n$ for every $n \geq N_\ast$, and hence
$h_n \leq \alpha^n\,h_0 \cdot M_\ast$ for a constant $M_\ast =
M_\ast(\alpha) < \infty$ absorbing the first $N_\ast$ levels. In
particular $h_n \to 0$ geometrically, and the induction closes
($h_n < \pi$ for all $n$).
\end{lemma}

\begin{proof}
Induction on $n$. Base: $h_0 = \pi/2 < \pi$
(Fact~\ref{fact:V0}), so arguments to $\mathrm{sph.mid}$ are
non-antipodal.

Step $n \to n + 1$: Assume $h_n < \pi$. Each level-$(n+1)$ edge
is either:
\begin{enumerate}
\item A midpoint-subdividing edge $\{u, m_{\{u, v\}}\}$ of a
  parent $\{u, v\} \in E(G_n^{D_4})$. By
  Lemma~\ref{lem:sph-mid-shrink}, its spherical length is
  $\tfrac{1}{2}\dsph(\iota_n(u), \iota_n(v)) \leq \tfrac{1}{2}
  h_n$.
\item A midpoint-to-centroid edge
  $\{m_{\{u, v\}}, c_F\}$ for a $2$-face $F = \{u, v, w\}$ (the
  $24$-cell's $2$-faces are triangles
  \cite[\S 22.2]{CoxeterRegularPolytopes}). The three vertices
  $p_1 := \iota_n(u), p_2 := \iota_n(v), p_3 := \iota_n(w) \in
  S^3$ have pairwise spherical distance at most $h_n$, and by
  induction (for $n \geq 1$) they have pairwise inner products
  $\langle p_i, p_j \rangle \geq \cos(h_n)$. Compute the spherical
  centroid and the midpoint-to-centroid distance explicitly.

  \emph{Base case $n = 0$}: at level $0$, $F$ is a $24$-cell
  $2$-face with vertices at pairwise inner product $0$
  (Fact~\ref{fact:V0}), i.e.\ all arcs equal $\pi/2 = h_0$.
  The centroid $c_F = (p_1 + p_2 + p_3)/\|p_1 + p_2 + p_3\|$
  satisfies $\|p_1 + p_2 + p_3\|^2 = 3 + 2 \cdot 0 \cdot 3 = 3$,
  so $c_F = (p_1 + p_2 + p_3)/\sqrt{3}$. The midpoint of edge
  $\{p_1, p_2\}$ is $m = (p_1 + p_2)/\|p_1 + p_2\| = (p_1 +
  p_2)/\sqrt{2}$ (using $\|p_1 + p_2\|^2 = 2 + 2 \cdot 0 = 2$).
  Inner product:
  \[
    \langle c_F, m \rangle = \frac{\langle p_1 + p_2 + p_3,
      p_1 + p_2\rangle}{\sqrt{3}\,\sqrt{2}}
    = \frac{\|p_1 + p_2\|^2 + \langle p_3, p_1 + p_2\rangle}
      {\sqrt{6}} = \frac{2 + 0}{\sqrt{6}} = \sqrt{2/3}.
  \]
  So $\dsph(c_F, m) = \arccos(\sqrt{2/3}) \approx 0.6155 <
  \pi/4 = h_0/2$, i.e.\ the midpoint-to-centroid edge at level
  $0 \to 1$ is strictly shorter than half of $h_0$. In
  particular $\dsph(c_F, m) < \tfrac{1}{2}\,h_0$.

  \emph{Inductive step $n \geq 1$}: the argument extends by
  continuity. For $p_1, p_2, p_3 \in S^3$ with pairwise inner
  products $\geq \cos h_n$, the centroid and midpoint computation
  above gives $\dsph(c_F, m) = F(\theta_{12}, \theta_{13},
  \theta_{23})$ as a smooth function of the pairwise angles
  $\theta_{ij} = \dsph(p_i, p_j) \leq h_n$. First-order Taylor
  expansion of the spherical formulas yields the linear-rate
  bound
  \[
    F(\theta_{12}, \theta_{13}, \theta_{23})
      \;\leq\; \frac{\max\theta_{ij}}{\sqrt{3}}
        + O\bigl(\max\theta_{ij}^3\bigr)
      \;\leq\; \frac{h_n}{\sqrt{3}} + O(h_n^3).
  \]
  At level $n = 0$ with all $\theta_{ij} = \pi/2 = h_0$, the
  exact value is $F = \arccos(\sqrt{2/3}) \approx 0.6155$; the
  bound $h_0/\sqrt{3} \approx 0.907$ is looser but still strictly
  less than $h_0 = \pi/2$. At level $n \geq 1$, the smaller
  $h_n$ makes the Taylor remainder negligible, and the bound
  $\dsph(c_F, m) \leq h_n/\sqrt{3} + O(h_n^3)$ holds uniformly.
\end{enumerate}
Both combinatorial classes give
$\ell_e^{(n+1)} \leq \max(\tfrac{1}{2},\,
\tfrac{1}{\sqrt{3}})\,h_n + O(h_n^3) = \tfrac{1}{\sqrt{3}}\,h_n
+ O(h_n^3)$. Fix any $\alpha \in (1/\sqrt{3}, 1)$, e.g.\
$\alpha = 0.6$. There exists a threshold level $N_\ast =
N_\ast(\alpha) < \infty$ (depending on how quickly the remainder
$O(h_n^3)$ is absorbed) such that $h_{n+1} \leq \alpha\,h_n$ for
all $n \geq N_\ast$; hence $h_n \leq \alpha^{n - N_\ast} h_{N_\ast}$
for all $n \geq N_\ast$. For the finitely many levels $n < N_\ast$
the bound $h_n \leq h_0$ is immediate from the subdivision rule,
so there exists an absolute constant $M_\ast = M_\ast(\alpha) \geq
1$ with $h_n \leq M_\ast \alpha^n h_0$ for \emph{all} $n \geq 0$;
in particular $h_n < \pi$ for all $n$, closing the induction.
(This establishes geometric decay with rate $\alpha$ up to the
absolute constant $M_\ast$; it does \emph{not} establish the
dyadic rate $h_n \leq 2^{-n} h_0$ with absolute constant $1$.)
\end{proof}

\begin{remark}[The constant $1/\sqrt{3}$]
\label{rmk:mid-cent-strict}
The constant $1/\sqrt{3}$ arises from the spherical centroid of
three vertices forming an isoceles triangle: in the linear
(flat-space) limit, the midpoint-to-centroid distance in an
equilateral triangle of side $s$ is $s/(2\sqrt{3})$, and for a
general triangle with max side $s$, it is at most $s/\sqrt{3}$
(by a geometric argument from the circumradius formula). The
base-case computation at $n = 0$, $\arccos(\sqrt{2/3}) \approx
0.6155$, is less than $h_0/\sqrt{3} \approx 0.907$, confirming
the bound with slack.
\end{remark}

\begin{hypothesis}[Uniform edge-length separation]
\label{fact:separation}
Let $e_n := \min_{\{u, v\} \in E(G_n^{D_4})} \dsph(\iota_n(u),
\iota_n(v))$. There exists a constant $c_e > 0$ (depending only on
the $24$-cell combinatorics and the Schl\"afli refinement rule)
such that $e_n \geq c_e h_n$ for all $n \geq 0$.
\end{hypothesis}

\begin{remark}[Status of Hypothesis~\ref{fact:separation}]
\label{rmk:separation-status}
At level $0$, $e_0 = h_0 = \pi/2$ so $e_0/h_0 = 1$, and the
hypothesis holds with $c_e \leq 1$. A full level-uniform
lower-bound argument for $e_n/h_n$ requires analysing the
spherical geometry of the refined triangles at each level,
specifically the minimum over all midpoint-to-centroid and
midpoint-subdividing edges; this is not carried out here. We
adopt Hypothesis~\ref{fact:separation} as an input to the
vertex-count / doubling results of \S\ref{sec:counting}; all
such results are conditional on it.
\end{remark}

\begin{hypothesis}[Spherical cells tile $S^3$]
\label{fact:tiling}
For every $n$, the spherical realization $\mathcal X_n^{\rm sph}$
(replacing each abstract cell of P3's level-$n$ complex by its
spherical convex hull via $\iota_n$) tiles $S^3$: the cells are
interior-disjoint with union $S^3$.
\end{hypothesis}

\begin{remark}[Status of Hypothesis~\ref{fact:tiling}]
\label{rmk:tiling-status}
At level $0$, $\mathcal X_0^{\rm sph}$ is the standard tessellation
of $S^3$ by the $24$ spherical octahedra of the $24$-cell
\cite[\S 22.2]{CoxeterRegularPolytopes}, coned over octahedron
barycenters to $48$ spherical tetrahedra; the tiling property
holds at $n = 0$. At higher levels, each Schl\"afli refinement
subdivides every level-$n$ spherical simplex into smaller
spherical simplices, and the level-$(n+1)$ cells are the union
of these refined subsimplices. The tiling property at level
$n+1$ requires two geometric inputs: (i) the refined subsimplices
of each parent are pairwise interior-disjoint with union equal
to the parent; (ii) refined subsimplices from adjacent parents
agree on shared boundary faces. Both are standard for Euclidean
barycentric subdivision; their rigorous transfer to spherical
subdivision via $\iota_n$ requires controlling the spherical
centroid positions consistently across parents, which this paper
does not carry out in full. We adopt
Hypothesis~\ref{fact:tiling} as a structural hypothesis of the
spherical realization; it is used in the $\varepsilon$-net
argument of \S\ref{sec:epsilon-net}. All theorems depending on
it are stated conditional on this hypothesis.
\end{remark}

%=====================================================================
\section{Intrinsic Graph Metric and $\varepsilon$-Net Theorem}
\label{sec:epsilon-net}
%=====================================================================

\begin{definition}[Edge length; two metrics on $V_n$]
\label{def:dn}
The spherical edge length is $\ell_e^{(n)} := \dsph(\iota_n(u),
\iota_n(v))$ for $e = \{u, v\} \in E(G_n^{D_4})$. We equip
$V_n$ with TWO metrics, for different purposes:
\begin{itemize}
\item The \emph{extrinsic metric} $d_n^{\rm ext}(u, v) :=
  \dsph(\iota_n(u), \iota_n(v))$, the pullback of $\dsph$ under
  $\iota_n$. This is the metric used for the (conditional on Hypothesis~\ref{fact:tiling}) GH
  convergence result (Theorem~\ref{thm:gh-convergence}).
\item The \emph{intrinsic graph metric} $d_n^{\rm int}$, the
  shortest-path distance on the edge-weighted graph
  $(G_n^{D_4}, \ell^{(n)})$. This is the metric used for the
  conditional Dirichlet-form / spectral-convergence program
  (\S\ref{sec:laplacian}--\S\ref{sec:spectral}).
\end{itemize}
By construction $d_n^{\rm ext}(u, v) \leq d_n^{\rm int}(u, v)$
(the shortest path in the graph is at least the straight-line
spherical distance, by triangle inequality). The
bi-Lipschitz comparison between $d_n^{\rm ext}$ and $d_n^{\rm
int}$ (i.e.\ an upper bound on their ratio) is a nontrivial
statement requiring path-reconstruction analysis; we present it
as part of the \emph{conditional program} of
\S\ref{sec:conditional-bilipschitz} below, under
Hypothesis~\ref{hyp:isotropy}. Unless stated otherwise, from
now on $d_n$ abbreviates $d_n^{\rm ext}$ in the base
statements.
\end{definition}

\begin{theorem}[$\varepsilon$-net, conditional on Hypothesis~\ref{fact:tiling}]
\label{thm:epsilon-net}
Assume Hypothesis~\ref{fact:tiling}. Then there exist constants
$C_1 > 0$ and $h_\ast \in (0, \pi/4)$ (depending only on the
$24$-cell combinatorial data) such that for every $n$ with
$h_n \leq h_\ast$: $\iota_n(V_n)$ is an $\varepsilon_n$-net in
$(S^3, \dsph)$ with $\varepsilon_n \leq C_1 h_n$. Moreover,
$\iota_n$ is an isometric embedding of $(V_n, d_n^{\rm ext})$
into $(S^3, \dsph)$ by Definition~\ref{def:dn}: $d_n^{\rm ext}(u, v) =
\dsph(\iota_n(u), \iota_n(v))$ for all $u, v \in V_n$ (this second
statement is by definition and does not itself require Hypothesis~\ref{fact:tiling}).
\end{theorem}

\begin{proof}
By Hypothesis~\ref{fact:tiling}, every $p \in S^3$ lies in some cell
of $\mathcal X_n^{\rm sph}$ with vertices in $\iota_n(V_n)$.
Spherical diameter of any cell is at most $c_{\rm cell}\,h_n$
(with $c_{\rm cell}$ a combinatorial constant $\leq 3$ for the
coned tetrahedral refinement). Pick the vertex of that cell
nearest to $p$: distance $\leq C_1 h_n$ with $C_1 :=
c_{\rm cell}$. The isometric-embedding statement is immediate
from Definition~\ref{def:dn} of the extrinsic metric.
\end{proof}

\begin{remark}[On the intrinsic-vs-extrinsic comparison]
\label{rmk:bilipschitz-conditional}
Earlier iterations of this paper attempted to prove a direct
bi-Lipschitz comparison between $d_n^{\rm int}$ and
$d_n^{\rm ext}$ unconditionally. The proof required a per-cell
geodesic-tracking bound that cannot be obtained without vertex-
figure isotropy at every vertex of every orbit; for a generic
spherical simplicial complex with bounded-degree vertices, the
multiplier in the bi-Lipschitz bound need not tend to $1$. We
defer the intrinsic-extrinsic comparison to
\S\ref{sec:conditional-bilipschitz} as part of the conditional
program under Hypothesis~\ref{hyp:isotropy}. For the
GH convergence of Theorem~\ref{thm:gh-convergence}
below, only the extrinsic metric $d_n^{\rm ext}$ is used.
\end{remark}

%=====================================================================
\section{Gromov--Hausdorff Convergence}
\label{sec:gh}
%=====================================================================

\begin{theorem}[GH convergence with explicit rate, conditional on Hypothesis~\ref{fact:tiling}]
\label{thm:gh-convergence}
\emph{Assume Hypothesis~\ref{fact:tiling}.}
Let $X_n := (V_n, d_n^{\rm ext})$ equipped with the extrinsic
metric of Definition~\ref{def:dn}, and $X_\infty := (S^3,
\dsph)$. For every $n$ with $h_n \leq h_\ast$,
\[
  \dGH(X_n, X_\infty) \;\leq\; \varepsilon_n \;\leq\; C_1\,h_n
    \;=\; O(h_n) \;=\; O(\alpha^n).
\]
\end{theorem}

\begin{proof}
$f_n := \iota_n \colon V_n \to S^3$ is an isometric embedding of
$(V_n, d_n^{\rm ext})$ into $(S^3, \dsph)$ by
Definition~\ref{def:dn} and Theorem~\ref{thm:epsilon-net}. By
Theorem~\ref{thm:epsilon-net}, $f_n(V_n) = \iota_n(V_n)$ is an
$\varepsilon_n$-net in $(S^3, \dsph)$ with $\varepsilon_n \leq
C_1 h_n$. An isometric $\varepsilon$-net embedding $f \colon
X \hookrightarrow Y$ gives $\dGH(X, Y) \leq \varepsilon$
\cite[Prop.~7.3.10]{BuragoBuragoIvanov}, hence $\dGH(X_n,
X_\infty) \leq \varepsilon_n \leq C_1 h_n$.
\end{proof}

%=====================================================================
\section{Intrinsic-vs-Extrinsic Comparison (Conditional Program)}
\label{sec:conditional-bilipschitz}
%=====================================================================

The intrinsic graph metric $d_n^{\rm int}$ (shortest-path
distance on $(G_n^{D_4}, \ell^{(n)})$) is needed for the
Dirichlet-form / spectral-convergence program of
\S\ref{sec:laplacian}--\S\ref{sec:spectral}, because the edge
weights of the Dirichlet form depend on the edge lengths. A
natural question is whether $d_n^{\rm int}$ agrees with
$d_n^{\rm ext}$ (modulo higher-order corrections) uniformly on
$V_n \times V_n$.

\begin{theorem}[Intrinsic-extrinsic comparison,
  conditional on Hypothesis~\ref{hyp:isotropy}]
\label{thm:intrinsic-extrinsic}
Under Hypothesis~\ref{hyp:isotropy} (uniform $8/3$ approximate
vertex-figure isotropy), there exists $C_2 > 0$ such that for
every $n$ with $h_n \leq h_\ast$ and every $u, v \in V_n$,
\[
  d_n^{\rm ext}(u, v) \;\leq\; d_n^{\rm int}(u, v)
    \;\leq\; d_n^{\rm ext}(u, v) + C_2\,h_n.
\]
\end{theorem}

\begin{remark}[Proof sketch and its gaps]
\label{rmk:bilipschitz-sketch}
The lower bound $d_n^{\rm ext} \leq d_n^{\rm int}$ is the
triangle inequality on $S^3$. The upper bound
$d_n^{\rm int} \leq d_n^{\rm ext} + C_2 h_n$ is the nontrivial
part; its proof requires a \emph{pearl-string} construction
along spherical geodesics (choose vertex approximants along
the geodesic at $\Theta(h_n)$-spaced intervals, connect
consecutive pearls by edges, bound per-edge excess via
second-order spherical geometry in a neighborhood). Each
pearl-connection step requires:
\begin{enumerate}
\item Forward-progress control: the greedy neighbor in the
  direction of $\iota_n(v)$ makes progress $h_n(1 + O(h_n^2))$
  per step. This requires the vertex-figure at every
  intermediate vertex to span enough directions, which is
  precisely Hypothesis~\ref{hyp:isotropy}.
\item Per-pearl excess bound: each step deviates from the
  geodesic by $O(h_n^2)$; summed over $L/h_n$ steps gives total
  excess $O(L h_n) \leq O(h_n)$.
\end{enumerate}
The detailed proof is a technical but standard greedy-step
argument in discrete metric geometry; we do not reproduce it
here. Without Hypothesis~\ref{hyp:isotropy}, the forward-
progress bound fails at midpoint / centroid orbits (where
vertex figures may not span enough directions), and the
multiplier in the bi-Lipschitz bound need not tend to $1$.
This is precisely why the upper bound is conditional.
\end{remark}

\begin{corollary}[GH convergence under the intrinsic metric,
  conditional]
\label{cor:gh-intrinsic}
Under Hypotheses~\ref{fact:tiling} and~\ref{hyp:isotropy},
$\dGH((V_n, d_n^{\rm int}), (S^3, \dsph)) = O(h_n)$ as well,
with the same rate as the extrinsic statement
of Theorem~\ref{thm:gh-convergence} (itself conditional on
Hypothesis~\ref{fact:tiling}).
\end{corollary}

\begin{proof}
Immediate from Theorem~\ref{thm:intrinsic-extrinsic} and the
extrinsic GH bound (Theorem~\ref{thm:gh-convergence}), plus
\cite[Cor.~7.3.28]{BuragoBuragoIvanov}.
\end{proof}

%=====================================================================
\section{Doubling, Packing, and Volume Comparison}
\label{sec:doubling}
%=====================================================================

\begin{definition}[Discrete and continuum measures]
\label{def:measures}
$m_n := |V_n|^{-1} \sum_{v \in V_n} \delta_v$, normalized counting
measure ($m_n(V_n) = 1$). $\mu_\infty$ is normalized Riemannian
volume on $(S^3, g_{\rm round})$, $\mu_\infty(S^3) = 1$.
\end{definition}

\begin{lemma}[Vertex count]
\label{lem:Vn-count}
There exist $c_V, C_V > 0$ with $c_V h_n^{-3} \leq |V_n| \leq
C_V h_n^{-3}$.
\end{lemma}

\begin{proof}
Upper: $\iota_n(V_n)$ are $c_e h_n$-separated on $S^3$
(Hypothesis~\ref{fact:separation}), so $|V_n| \leq
\vol_{S^3}(S^3)/\vol_{S^3}(B(p, c_e h_n/2)) \leq C_V h_n^{-3}$ by
round-sphere volume estimate \cite[Ch.~III]{BerardSpectralGeometry}.
Lower: $\iota_n(V_n)$ is a $C_1 h_n$-net
(Theorem~\ref{thm:epsilon-net}), so $|V_n| \geq \vol(S^3)/
\vol(B(p, C_1 h_n)) \geq c_V h_n^{-3}$.
\end{proof}

\begin{proposition}[Doubling and packing, conditional on Hypotheses~\ref{fact:tiling} and~\ref{fact:separation}, extrinsic]
\label{prop:doubling}
\emph{Assume Hypotheses~\ref{fact:tiling} and~\ref{fact:separation}.}
Let $r_{\min} := 2 C_1 h_n$. There exist $c_0, C_0, C_D > 0$
such that for every $n$ with $h_n \leq h_\ast$, every $v \in
V_n$, and every $r \in [r_{\min}, \pi/4]$, using the
\emph{extrinsic}-metric ball $B_{d_n^{\rm ext}}(v, r) := \{u
\in V_n : \dsph(\iota_n(u), \iota_n(v)) \leq r\}$:
\begin{enumerate}
\item[\textup{(i)}] $c_0 r^3 \leq m_n(B_{d_n^{\rm ext}}(v, r))
  \cdot (2\pi^2) \leq C_0 r^3$.
\item[\textup{(ii)}] $m_n(B_{d_n^{\rm ext}}(v, 2r)) \leq
  C_D\,m_n(B_{d_n^{\rm ext}}(v, r))$.
\item[\textup{(iii)}] Any $d_n^{\rm ext}$-$r$-separated subset
  of $V_n$ has $m_n$-measure at most $C_P (h_n/r)^3$ with $C_P
  := 8/c_0$.
\end{enumerate}
\end{proposition}

\begin{proof}
Under the extrinsic metric, $\iota_n$ is isometric
(Theorem~\ref{thm:epsilon-net}), so $\iota_n(B_{d_n^{\rm ext}}(v,
r)) = B_{\dsph}(\iota_n(v), r) \cap \iota_n(V_n)$ exactly (no
additive distortion needed).

Count of $\iota_n(V_n)$-points in a round ball of radius $s \in
[c_e h_n, \pi/4]$ is $\Theta(s^3 / h_n^3)$: the covering
direction uses Theorem~\ref{thm:epsilon-net} ($\iota_n(V_n)$
is a $C_1 h_n$-net), and the separation direction uses
Hypothesis~\ref{fact:separation} ($c_e h_n$-separated). Both give
$s^3 / h_n^3$ up to explicit constants. Normalizing by $|V_n| =
\Theta(h_n^{-3})$ (Lemma~\ref{lem:Vn-count}) gives
$m_n(B) = \Theta(s^3)$, i.e.\ (i) holds for $r$ in the adjusted
admissible range. Doubling (ii) is immediate from $c_0 (2r)^3
\leq C_D c_0 r^3$ with $C_D = 8 C_0/c_0$. Packing (iii) from
$c_0 r^3 \leq |V_n|\,m_n(B)$ applied to each
ball in the disjoint family.
\end{proof}

%=====================================================================
\section{Vertex-Figure Isotropy: The Structural Constant $8/3$}
\label{sec:isotropy}
%=====================================================================

\begin{notation}[Tangential projection $\pi_v$]
\label{not:pi-v}
For $v \in V_n$, $T_{\iota_n(v)} S^3 := \{x \in \R^4 : \langle
x, \iota_n(v)\rangle = 0\}$, $\pi_v(x) := x - \langle x,
\iota_n(v)\rangle\,\iota_n(v)$.
\end{notation}

\begin{definition}[Vertex-figure tensor (unweighted and weighted)]
\label{def:T-v}
The \emph{unweighted vertex-figure tensor} at $v \in V_n$ is
\[
  T_v^{(n)} \;:=\; \sum_{u : \{u, v\} \in E(G_n^{D_4})}
    \frac{\pi_v(\iota_n(u) - \iota_n(v))
    \otimes \pi_v(\iota_n(u) - \iota_n(v))}
    {\|\pi_v(\iota_n(u) - \iota_n(v))\|^2}
  \;\in\; \Sym^2(T_{\iota_n(v)} S^3).
\]
The \emph{weighted vertex-figure tensor}, more directly relevant
to the Dirichlet-form consistency of \S\ref{sec:laplacian}, is
\[
  T_v^{(n),\,\rm w} \;:=\; \frac{1}{h_n^2}
    \sum_{u : \{u, v\} \in E(G_n^{D_4})}
    \pi_v(\iota_n(u) - \iota_n(v))
    \otimes \pi_v(\iota_n(u) - \iota_n(v))
  \;=\; \frac{1}{h_n^2}\sum_{u \sim v}
    (\ell_e^{(n)})^2 \cdot \widehat{\tau}_e \otimes \widehat{\tau}_e,
\]
where $\widehat{\tau}_e := \pi_v(\iota_n(u) - \iota_n(v))/
\|\pi_v(\iota_n(u) - \iota_n(v))\|$ is the unit tangent direction
along edge $e = \{u, v\}$ (after rotating the spherical arc
identification: the unit tangent lies in $T_{\iota_n(v)} S^3$).
The two tensors agree up to diagonal rescaling by
$(\ell_e^{(n)})^2/h_n^2 \in (c_e^2, 1]$.
\end{definition}

\begin{proposition}[Exact isotropy at $V_0$ vertices, every level]
\label{prop:v0-at-level-n}
For every $v \in V_0$ and every $n \geq 0$ (viewing $v$ as a
vertex of $V_n$ via the vertex inclusion $V_0 \subset V_n$),
\[
  T_v^{(n)} \;=\; \tfrac{8}{3}\,I_{T_{\iota_n(v)} S^3}.
\]
\end{proposition}

\begin{proof}
By $W(D_4)$-transitivity on $V_{24}$ (Notation~\ref{not:WD4}), it
suffices to verify the identity at one representative, say $v =
e_1$. The level-$n$ neighbors of $v = e_1$ in $G_n^{D_4}$ are
either (at level $n = 0$) the $8$ half-type vertices
$u_\varepsilon := \tfrac{1}{2}(1, \varepsilon_2, \varepsilon_3,
\varepsilon_4)$, $\varepsilon \in \{\pm 1\}^3$, or (at level $n
\geq 1$) the $8$ spherical midpoints of the $8$ original edges
at $v$ under iterated refinement (each such midpoint lies along
the great-circle arc from $v$ to $u_\varepsilon$ at spherical
distance $2^{-n} \cdot \pi/2$, by
Lemma~\ref{lem:sph-mid-shrink}).

\emph{Case $n = 0$.} For each $\varepsilon$:
$u_\varepsilon - v = \tfrac{1}{2}(-1, \varepsilon_2, \varepsilon_3,
\varepsilon_4)$; $\langle u_\varepsilon - v, v\rangle = -1/2$;
$\pi_v(u_\varepsilon - v) = (0, \varepsilon_2/2, \varepsilon_3/2,
\varepsilon_4/2)$; $\|\pi_v(\cdot)\|^2 = 3/4$. Sum of outer
products in coordinates $(e_2, e_3, e_4)$:
\[
  \sum_{\varepsilon \in \{\pm 1\}^3}
    \pi_v(u_\varepsilon - v) \otimes \pi_v(u_\varepsilon - v)
    \;=\; \tfrac{1}{4} \sum_\varepsilon \varepsilon_i
       \varepsilon_j \delta_{(i, j)}
    \;=\; \tfrac{1}{4} \cdot 8 \cdot I_3 = 2\,I_3
\]
(off-diagonal entries vanish by parity). Dividing by $\|\pi_v\|^2
= 3/4$ per term: $T_v^{(0)} = (2\,I_3)/(3/4) = (8/3)\,I_3 =
(8/3)\,I_{T_v S^3}$.

\emph{Case $n \geq 1$.} The $8$ level-$n$ midpoints are
$m_\varepsilon^{(n)}$, lying on the great-circle arcs from $v$ to
$u_\varepsilon$ at spherical distance $\theta_n := 2^{-n} \cdot
\pi/2$ from $v$. Explicitly:
$m_\varepsilon^{(n)} = \cos(\theta_n)\,v +
\sin(\theta_n)\,\tau_\varepsilon$, where $\tau_\varepsilon$ is
the unit tangent at $v$ toward $u_\varepsilon$:
$\tau_\varepsilon = \pi_v(u_\varepsilon - v)/\|\pi_v(u_\varepsilon
- v)\| = (0, \varepsilon_2, \varepsilon_3,
\varepsilon_4)/\sqrt{3}$.

Thus
\[
  m_\varepsilon^{(n)} - v = (\cos\theta_n - 1)\,v +
    \sin\theta_n\,\tau_\varepsilon,
\]
$\pi_v(m_\varepsilon^{(n)} - v) = \sin\theta_n\,\tau_\varepsilon
= \sin\theta_n\,(0, \varepsilon_2, \varepsilon_3,
\varepsilon_4)/\sqrt{3}$. $\|\pi_v(\cdot)\|^2 = \sin^2\theta_n$.
Sum of outer products:
\[
  \sum_\varepsilon \pi_v \otimes \pi_v
    \;=\; \sin^2\theta_n \cdot \tfrac{1}{3}
       \sum_\varepsilon \varepsilon_i\varepsilon_j
    \;=\; \sin^2\theta_n \cdot \tfrac{8}{3}\,I_3.
\]
Divided by $\sin^2\theta_n$ per term:
$T_v^{(n)} = (8/3)\,I_3 = (8/3)\,I_{T_v S^3}$.

The $8/3$ is identical at every level: the tangent directions
$\tau_\varepsilon$ remain the same (midpoints lie along the
\emph{same} great-circle arcs), only their radial distance
shrinks, and the normalization by $\|\pi_v\|^2$ cancels the
shrinkage.
\end{proof}

\begin{remark}[Origin of the $8/3$]
\label{rmk:origin-83}
The constant $8/3$ is the outer-product sum normalization of the
cube in $\R^3$: the $8$ cube vertices $(\varepsilon_2,
\varepsilon_3, \varepsilon_4)/\sqrt{3}$ on the unit sphere in
$\R^3$ sum as
$\sum_\varepsilon (\varepsilon/\sqrt{3}) \otimes
(\varepsilon/\sqrt{3}) = (8/3)\,I_3$. This is an elementary
identity; its non-trivial role here is that the $D_4$ Schl\"{a}fli
vertex figure at every $V_0$ vertex \emph{is} such a cube,
preserved at every level by $W(D_4)$-equivariance.
\end{remark}

\begin{hypothesis}[Uniform $8/3$ approximate weighted isotropy]
\label{hyp:isotropy}
For every $n \geq 0$ and every $v \in V_n$ (including midpoint
and centroid orbits), the \emph{weighted} vertex-figure tensor
satisfies
\[
  T_v^{(n),\,\rm w} \;=\; \tfrac{8}{3}\,I_{T_{\iota_n(v)} S^3}
    + E_v^{(n)},
  \qquad \|E_v^{(n)}\|_{\rm op} \leq C_{\rm iso}\,h_n^2,
\]
for a single constant $C_{\rm iso}$ independent of $n$ and $v$.
The constant $8/3$ matches the $V_0$ orbit value (where
$\ell_e^{(n)} = h_n$ exactly and the unweighted and weighted
tensors coincide). The hypothesis is about the \emph{weighted}
sum $\sum (\ell_e^{(n)}/h_n)^2\,\widehat{\tau}_e \otimes
\widehat{\tau}_e$ of Definition~\ref{def:T-v}, which is the
quantity that appears directly in the Dirichlet-form
consistency analysis of \S\ref{sec:laplacian}.
\end{hypothesis}

\begin{remark}[Status of Hypothesis~\ref{hyp:isotropy}]
\label{rmk:hyp-status}
Hypothesis~\ref{hyp:isotropy} is PROVED unconditionally for the
$V_0$ orbit at every level (Proposition~\ref{prop:v0-at-level-n},
with $E_v^{(n)} = 0$ and the constant $8/3$ exact).
For midpoint and centroid
orbits in $V_n \setminus V_0$ at $n \geq 1$, it is plausible by
$W(D_4)$-equivariance + smooth-limit Taylor expansion but not
proved here; completion is the subject of a follow-up paper.
All results in \S\ref{sec:laplacian}--\S\ref{sec:spectral}
depend on this hypothesis, stated explicitly.
\end{remark}

%=====================================================================
\section{Graph Laplacian and Dirichlet-Form Consistency}
\label{sec:laplacian}
%=====================================================================

\begin{definition}[Graph Laplacian $\Delta_n$]
\label{def:Delta-n}
Set $c_\Delta := 3/4$, the normalization constant derived in
Step 3 of Theorem~\ref{thm:form-consistency} below: $c_\Delta =
(8/3)^{-1} \cdot 2$ $= 3/4$, balancing the weighted isotropy
constant $8/3$ of Hypothesis~\ref{hyp:isotropy} with the
factor-of-$2$ edge double-counting in the per-vertex sum versus
edge sum. The graph Laplacian on $L^2(V_n, m_n)$
is
\[
  (\Delta_n f)(v)
    \;:=\; \frac{c_\Delta\,|V_n|}{h_n^2}
    \sum_{u : \{u, v\} \in E(G_n^{D_4})} (f(v) - f(u)).
\]
The Dirichlet form is $\mathcal E_n(f, g) := \langle \Delta_n f,
g \rangle_{L^2(V_n, m_n)}$, giving
\[
  \mathcal E_n(f) \;=\; \frac{c_\Delta}{h_n^2 |V_n|^{-1}}
    \sum_{\{u, v\}} (f(v) - f(u))^2 / |V_n|
    \;=\; \frac{c_\Delta}{h_n^2}\,\langle \text{edge sum}
    \rangle_{|V_n|^{-1}}.
\]
The isotropy constant $8/3$ and the edge-double-count factor
$1/2$ together motivate $c_\Delta = 3/4$, so that the Dirichlet
form converges to the continuum energy at leading order.
\end{definition}

\begin{remark}[Edge-length dependence addressed]
\label{rmk:edge-length}
Though $\Delta_n$ as defined above is formally an unweighted
adjacency Laplacian times scalars, the \emph{normalization}
$c_\Delta$ and the rescaling $1/h_n^2$ together embed the
edge-length structure: edge lengths are all $\Theta(h_n)$ by
Hypothesis~\ref{fact:separation}, and the continuum-limit computation
(Theorem~\ref{thm:form-consistency}) uses $(f(u) - f(v))^2 =
(\ell_e^{(n)})^2\,(\nabla_\tau f)^2 + O(\ell_e^{(n), 3})$, with
the $\ell^2$ factors appearing explicitly in the Taylor expansion.
The unweighted adjacency form is the correct leading-order
approximation under Hypothesis~\ref{hyp:isotropy}.
\end{remark}

\begin{definition}[Transport $R_n, I_n$]
\label{def:transport}
$R_n \phi(v) := \fint_{B_{\dsph}(\iota_n(v), C_1 h_n)}
\phi \, d\mu_\infty$ (ball average);
$I_n f(p) := \sum_v f(v) \psi_v(p)$ for a smooth partition of
unity $\{\psi_v\}$ subordinate to $\{B_{\dsph}(\iota_n(v), 2
C_1 h_n)\}$, with bounded overlap $N_{\rm ov}$ (from
Proposition~\ref{prop:doubling}) and $\|\nabla
\psi_v\|_\infty \leq C_\psi / h_n$.
\end{definition}

\begin{proposition}[Transport boundedness]
\label{prop:transport-bounded}
$\|R_n\|_{\rm op}$ and $\|I_n\|_{\rm op}$ are uniformly bounded
in $n$, with constants depending only on the $24$-cell data and
the partition-of-unity choice.
\end{proposition}

\begin{proof}[Sketch]
$R_n$: Jensen's inequality on the average + doubling
(Proposition~\ref{prop:doubling}) + bounded overlap gives
$\|R_n \phi\|_{L^2(m_n)}^2 \leq (N_{\rm ov}/c_0)\,\|\phi\|_{L^2(\mu_\infty)}^2$
(where $c_0$ is the doubling lower-volume constant from
Proposition~\ref{prop:doubling})
after $|V_n|\,h_n^3 = \Theta(1)$ normalization from
Lemma~\ref{lem:Vn-count}.

$I_n$: Cauchy--Schwarz + partition summation gives $\|I_n
f\|_{L^2(\mu_\infty)}^2 \leq N_{\rm ov}\,\sum_v |f(v)|^2
\mu_\infty(\supp \psi_v) \leq C h_n^3 |V_n|\,\|f\|_{L^2(m_n)}^2
= O(1)\,\|f\|_{L^2(m_n)}^2$ by Lemma~\ref{lem:Vn-count}.
\end{proof}

\begin{theorem}[Dirichlet-form consistency, conditional]
\label{thm:form-consistency}
Under Hypothesis~\ref{hyp:isotropy}, there exists $C_4 > 0$
independent of $n, f$ such that for every $f \in C^2(S^3)$ and
every $n$ with $h_n \leq h_\ast$,
\[
  \bigl|\mathcal E_n(R_n f) - \mathcal E_\infty(f)\bigr|
    \;\leq\; C_4\,h_n\,\|f\|_{C^2(S^3)}^2.
\]
The estimate is stated in the $C^2$-norm, not the $H^2$-norm;
see Remark~\ref{rmk:c2-norm} for why an $H^2$-rate extension
does not follow by density alone in dimension $3$.
\end{theorem}

\begin{proof}
We prove the estimate for $f \in C^2(S^3)$, where pointwise
second derivatives are well-defined.

\emph{Step 1: Pointwise replacement of $R_n f$.} For $f \in
C^2(S^3)$, the ball average $R_n f$ satisfies, at each
vertex $v$:
\[
  (R_n f)(v) \;=\; f(\iota_n(v)) + \tfrac{1}{2}
    \langle \nabla^2 f(\iota_n(v)),\, \mathrm{Var}_{B_v}\rangle
    + O(h_n^3 \|f\|_{C^2}),
\]
where $\mathrm{Var}_{B_v}$ is the variance tensor of uniform
distribution on $B_{\dsph}(\iota_n(v), C_1 h_n)$. On the round
sphere $(S^3, g_{\rm round})$, this variance is
$O(h_n^2)\,I_{T_v S^3}$. Thus $(R_n f)(v) = f(\iota_n(v)) +
O(h_n^2\,\|f\|_{C^2})$. Hence
$(R_n f)(v) - (R_n f)(u) = f(\iota_n(v)) - f(\iota_n(u)) +
O(h_n^2\,\|f\|_{C^2})$. Squaring and summing over
$|E(G_n^{D_4})| = O(|V_n|) = O(h_n^{-3})$ edges, and normalizing
by the $c_\Delta/(h_n^2 |V_n|)$ factor of
Definition~\ref{def:Delta-n}: the $R_n$-replacement error
contributes
\[
  c_\Delta\,h_n^{-2} |V_n|^{-1} \cdot O(|V_n|) \cdot
    O(h_n^2 \|f\|_{C^2} \cdot h_n \|f\|_{C^1})
  \;=\; O(h_n \|f\|_{C^2}^2)
\]
to $|\mathcal E_n(R_n f) - \mathcal E_n(f \circ \iota_n)|$,
where we used the pointwise product $|(R_n f - f \circ \iota_n)
\cdot (f \circ \iota_n)| \leq h_n^2 \|f\|_{C^2} \cdot h_n
\|f\|_{C^1}$ per edge.

\emph{Step 2: Taylor expansion of $f \circ \iota_n$ along
spherical edges.} For $f \in C^2(S^3)$ and an edge $e = \{u, v\}$
with $p := \iota_n(v)$, $q := \iota_n(u)$, let
$\gamma_e \colon [0, \ell_e^{(n)}] \to S^3$ be the unit-speed
geodesic from $p$ to $q$ (exists uniquely since $\ell_e^{(n)}
\leq h_n < \pi$). By the fundamental theorem of calculus,
\[
  f(q) - f(p) \;=\; \int_0^{\ell_e^{(n)}} \nabla_\tau f(\gamma_e(s))
    \, ds \;=\; \ell_e^{(n)} \cdot \nabla_\tau f(p) +
    O\bigl((\ell_e^{(n)})^2 \,\|\nabla^2 f\|_{C^0}\bigr),
\]
where $\tau := \gamma_e'(0)$ is the initial tangent. Squaring,
\[
  (f(q) - f(p))^2 \;=\; (\ell_e^{(n)})^2\,(\nabla_\tau f(p))^2
    + O\bigl((\ell_e^{(n)})^3\,\|f\|_{C^1}\,\|f\|_{C^2}\bigr).
\]

\emph{Step 3: Weighted isotropy summation and derivation of
$c_\Delta$.} Summing the squared edge differences at a vertex
$v$, using the \emph{weighted} isotropy tensor
$T_v^{(n),\,\rm w}$ of Definition~\ref{def:T-v} and
Hypothesis~\ref{hyp:isotropy}:
\begin{align*}
  \sum_{u \sim v}(\ell_e^{(n)})^2 (\nabla_\tau f(v))^2
    &\;=\; \sum_{u \sim v}(\ell_e^{(n)})^2\,
       \widehat{\tau}_e^\top (\nabla f \otimes \nabla f)\,\widehat{\tau}_e
    \;=\; \mathrm{tr}\bigl((\nabla f \otimes \nabla f)
       \cdot h_n^2\,T_v^{(n),\,\rm w}\bigr) \\
    &\;=\; h_n^2 \cdot \mathrm{tr}\bigl((\nabla f \otimes \nabla f)
       \cdot (\tfrac{8}{3}\,I + E_v^{(n)})\bigr)
    \;=\; h_n^2\,\bigl(\tfrac{8}{3}\,|\nabla f(v)|^2 +
       O(h_n^2\,\|f\|_{C^1}^2)\bigr).
\end{align*}
Summing over $v \in V_n$:
$\sum_v \sum_{u \sim v} (\ell_e^{(n)})^2 (\nabla_\tau f)^2 =
h_n^2 \cdot (8/3) \sum_v |\nabla f(v)|^2 + O(h_n^4 |V_n|
\|f\|_{C^1}^2)$.

Each undirected edge $\{u, v\} \in E(G_n^{D_4})$ contributes
twice to the per-vertex double sum (once from each endpoint),
so
\[
  \sum_{\{u, v\} \in E}(\ell_e^{(n)})^2 (\nabla_\tau f)^2
    \;=\; \tfrac{1}{2} \cdot h_n^2 \cdot \tfrac{8}{3}\,
       \sum_v |\nabla f(v)|^2 + O(h_n^4 |V_n| \|f\|_{C^1}^2)
    \;=\; \tfrac{4}{3}\,h_n^2\,\sum_v |\nabla f(v)|^2 + \ldots
\]

Substituting into $\mathcal E_n(f \circ \iota_n) = (c_\Delta/(h_n^2
|V_n|)) \sum_{\{u, v\}} (f(\iota_n(u)) - f(\iota_n(v)))^2$ and
using $(f(\iota_n(u)) - f(\iota_n(v)))^2 = (\ell_e^{(n)})^2
(\nabla_\tau f)^2 + O((\ell_e^{(n)})^3 \|f\|_{C^2})$ from Step 2:
\[
  \mathcal E_n(f \circ \iota_n) \;=\; \frac{c_\Delta\,(4/3)}{|V_n|}
    \sum_{v \in V_n} |\nabla f(v)|^2 + O(h_n\,\|f\|_{C^2}^2).
\]
For the discrete average $(1/|V_n|)\sum_v |\nabla f|^2(\iota_n(v))
\to \mathcal E_\infty(f)$ (Step 4 Riemann-sum quadrature below),
the continuum-limit prefactor of $\mathcal E_n$ must equal $1$;
i.e.\ $c_\Delta \cdot (4/3) = 1$, giving
\[
  \boxed{c_\Delta = 3/4.}
\]
This is the calibration choice in
Definition~\ref{def:Delta-n}.

\emph{Step 4: Riemann-sum quadrature on the
$\varepsilon_n$-net.} By Lemma~\ref{lem:quadrature} below
(which is a standard manifold-quadrature estimate under our
proved doubling Prop.~\ref{prop:doubling} and vertex-count
Lemma~\ref{lem:Vn-count}), for every $g \in C^1(S^3)$:
\[
  \biggl|\,\frac{1}{|V_n|}\sum_{v \in V_n} g(\iota_n(v))
    - \int_{S^3} g \, d\mu_\infty\,\biggr|
    \;\leq\; C_Q\,h_n\,\|g\|_{C^1(S^3)}.
\]
Applied to $g := |\nabla f|^2 \in C^1(S^3)$ (since
$f \in C^2$) with $\|g\|_{C^1} \leq 2 \|f\|_{C^1} \|f\|_{C^2}
\leq \|f\|_{C^2}^2$:
$(1/|V_n|)\sum_v |\nabla f|^2(\iota_n(v)) = \int_{S^3}
|\nabla f|^2 d\mu_\infty + O(h_n\,\|f\|_{C^2}^2)$.

Combining Steps 1--4 gives
$|\mathcal E_n(R_n f) - \mathcal E_\infty(f)| \leq C_4\,h_n\,
\|f\|_{C^2}^2$ for $f \in C^2(S^3)$, with explicit $C_4$
determined by $C_Q$, $C_{\rm iso}$, $c_0, C_0, C_D$, and $c_e$.
\end{proof}

\begin{lemma}[Riemann-sum quadrature on the $\varepsilon_n$-net]
\label{lem:quadrature}
There exists $C_Q > 0$ such that for every $n$ with $h_n \leq
h_\ast$ and every $g \in C^1(S^3)$,
\[
  \biggl|\,\frac{1}{|V_n|}\sum_{v \in V_n} g(\iota_n(v))
    - \int_{S^3} g \, d\mu_\infty\,\biggr|
    \;\leq\; C_Q\,h_n\,\|g\|_{C^1(S^3)}.
\]
\end{lemma}

\begin{proof}[Sketch]
Standard midpoint-rule quadrature on a compact Riemannian
manifold, using: (i) the $\varepsilon_n$-net structure
(Theorem~\ref{thm:epsilon-net}) gives a Voronoi partition of
$S^3$ into cells of diameter $\leq 2 C_1 h_n$; (ii) each cell
has $\mu_\infty$-volume equal to $1/|V_n|$ up to $O(h_n^4/|V_n|)$
error, by the doubling estimate
(Proposition~\ref{prop:doubling}) and the uniform vertex-count
asymptotic (Lemma~\ref{lem:Vn-count}); (iii) on each cell, $g$
is approximated by its value at the vertex with error $O(h_n
\|g\|_{C^1})$. Summing over cells gives the claimed bound. See
\cite[Ch.~4]{BrennerScottFEM} for the analogous planar
finite-element quadrature argument.
\end{proof}

\begin{remark}[On the $C^2$-norm in the consistency rate]
\label{rmk:c2-norm}
The estimate $|\mathcal E_n(R_n f) - \mathcal E_\infty(f)| \leq
C_4 h_n \|f\|_{C^2}^2$ is stated in the $C^2$-norm rather than
$H^2$-norm. In dimension $3$, $H^2(S^3) \hookrightarrow C^0(S^3)$
(Sobolev embedding) but \emph{not} into $C^2$; extending the
rate from $C^2$ to $H^2$ by density alone would lose the rate.
For the spectral convergence application of
Theorem~\ref{thm:spectral-convergence}, we apply
Theorem~\ref{thm:form-consistency} to specific smooth
eigenfunctions $\phi_j \in C^\infty(S^3)$ of $\Delta_\infty$
(spherical harmonics), whose $C^2$-norms are bounded by explicit
polynomial functions of the eigenvalue; the $C^2$-norm is the
natural regularity for this pointwise-Taylor-based argument.
\end{remark}

% Lemma H2-density removed in iter-6: the C^2 → H^2 density
% extension does not preserve the rate in dimension 3 (H^2 ↛ C^2
% is not a Sobolev embedding). The C^2 rate is the honest
% statement of Theorem~\ref{thm:form-consistency}, sufficient for
% the spectral application to smooth eigenfunctions.

%=====================================================================
\section{Spectral Convergence via Courant--Fischer (conditional)}
\label{sec:spectral}
%=====================================================================

\begin{theorem}[Eigenvalue convergence, conditional]
\label{thm:spectral-convergence}
Under Hypothesis~\ref{hyp:isotropy}, for every $k \geq 1$,
$\lambda_k^{(n)} \to \lambda_k^\infty$ as $n \to \infty$ (with
multiplicity), where $\lambda_k^{(n)}$ and $\lambda_k^\infty$
are the $k$th ordered eigenvalues of $\Delta_n$ and
$\Delta_\infty := -\Delta_{S^3}$.
\end{theorem}

\begin{proof}
Courant--Fischer min-max \cite[\S VIII.8]{Kato}:
\[
  \lambda_k^{(n)} = \inf_{\substack{F \subset L^2(V_n, m_n)\\
       \dim F = k}} \sup_{f \in F, \|f\| = 1}
       \mathcal E_n(f),
  \qquad \lambda_k^\infty = \inf_{F \subset L^2(S^3, \mu_\infty)}
       \ldots
\]

\emph{Upper bound $\limsup \lambda_k^{(n)} \leq \lambda_k^\infty$.}
Let $F_k^\infty := \mathrm{span}(\phi_1, \ldots, \phi_k) \subset
C^\infty(S^3)$ be the span of the first $k$ $\Delta_\infty$-
eigenfunctions (spherical harmonics, $\|\phi_j\|_{L^2} = 1$).
The eigenfunctions are smooth with explicit $C^2$-bounds:
$\sup_{j \leq k} \|\phi_j\|_{C^2(S^3)} =: C_k^\infty < \infty$
(for spherical harmonics of degree $\ell$, $\|\phi_j\|_{C^2} =
O(\ell^2)$ by standard spherical-harmonic estimates
\cite[Ch.~III]{BerardSpectralGeometry}).

\emph{Rank preservation for $R_n$.} We claim that for
$h_n \leq h_\ast^{(k)}$ (a threshold depending on $k$),
$R_n \colon F_k^\infty \to L^2(V_n, m_n)$ is injective. For
$f \in F_k^\infty$ smooth, the local-average approximation
$(R_n \phi)(v) = \phi(\iota_n(v)) + O(h_n^2 \|\phi\|_{C^2})$
(see Step 1 of Theorem~\ref{thm:form-consistency}) combined
with Lemma~\ref{lem:quadrature} applied to $|\phi|^2$ gives
$\|R_n f\|_{L^2(V_n, m_n)}^2 = (1/|V_n|)\sum_v |f(\iota_n(v))|^2
+ O(h_n^2 \|f\|_{C^2}^2) = \|f\|_{L^2(S^3)}^2 + O(h_n \|f\|_{C^2}^2)$.
For $f \in F_k^\infty$ smooth with $\|f\|_{L^2} = 1$ and
$\|f\|_{C^2} \leq C_k^\infty$, $\|R_n f\|_{L^2(m_n)}^2 \geq 1
- O(h_n (C_k^\infty)^2) \geq 1/2$ for $h_n$ small, so $R_n f
\neq 0$. Injectivity on $F_k^\infty$ follows by linearity.

Let $F_k^{(n)} := R_n(F_k^\infty) \subset L^2(V_n, m_n)$ of
dimension $k$. For $\tilde f = R_n f$ with $f \in F_k^\infty$,
by Theorem~\ref{thm:form-consistency} applied in the $C^2$-norm,
\[
  \mathcal E_n(\tilde f) \;\leq\; \mathcal E_\infty(f) + C_4 h_n
    \|f\|_{C^2}^2 \;\leq\; \lambda_k^\infty + C_4 h_n
    (C_k^\infty)^2.
\]
Min-max gives $\lambda_k^{(n)} \leq \sup_{\tilde f \in F_k^{(n)},
\|\tilde f\|=1} \mathcal E_n(\tilde f) \leq \lambda_k^\infty +
O(h_n)$, so $\limsup \lambda_k^{(n)} \leq \lambda_k^\infty$.

\emph{Lower bound $\liminf \lambda_k^{(n)} \geq \lambda_k^\infty$.}
Symmetric argument using the dual Dirichlet-form-consistency
estimate of Lemma~\ref{lem:dual-form-consistency} below.

Let $F_k^{(n)} := \mathrm{span}(\phi_1^{(n)}, \ldots,
\phi_k^{(n)}) \subset L^2(V_n, m_n)$ be spanned by the first $k$
$\Delta_n$-eigenfunctions with $\|\phi_j^{(n)}\|_{L^2(m_n)} = 1$.
$I_n(F_k^{(n)}) \subset L^2(S^3, \mu_\infty)$ has dimension at
most $k$, and is in $H^1(S^3)$ (since $I_n$-images are smooth by
the partition-of-unity construction, with $H^1$-norms controlled
by the discrete Dirichlet form).

We claim $I_n$ preserves rank on $F_k^{(n)}$ for $h_n$ small
enough: indeed, $R_n I_n f = f + O(h_n)\|f\|$ for $f \in F_k^{(n)}$
(by a symmetric argument to rank preservation of $R_n$, using
$H^1$-uniform bounds on the finite-dim eigenspace of $\Delta_n$,
a standard discrete elliptic regularity).

For $\phi \in I_n(F_k^{(n)})$ with $\|\phi\|_{L^2(S^3)} = 1$ and
pre-image $f \in F_k^{(n)}$: by
Lemma~\ref{lem:dual-form-consistency},
\[
  \mathcal E_\infty(\phi) \;\leq\; \mathcal E_n(f) + C_5 h_n
    \|f\|^2_{\mathcal E_n + L^2}
\]
(where $\|f\|^2_{\mathcal E_n + L^2} := \mathcal E_n(f) +
\|f\|_{L^2(m_n)}^2$ is the graph energy norm). For
$f = \sum_j c_j \phi_j^{(n)}$ with $\sum_j c_j^2 = 1$,
$\mathcal E_n(f) = \sum_j c_j^2 \lambda_j^{(n)} \leq
\lambda_k^{(n)}$, so $\mathcal E_\infty(\phi) \leq \lambda_k^{(n)}
+ O(h_n)(\lambda_k^{(n)} + 1)$. Min-max gives $\lambda_k^\infty
\leq \sup_\phi \mathcal E_\infty(\phi) \leq \lambda_k^{(n)} +
O(h_n)$, so $\liminf \lambda_k^{(n)} \geq \lambda_k^\infty$.

Combining upper and lower bounds: $\lambda_k^{(n)} \to
\lambda_k^\infty$.
\end{proof}

\begin{lemma}[Dual Dirichlet-form consistency for $I_n$]
\label{lem:dual-form-consistency}
Under Hypothesis~\ref{hyp:isotropy}, there exists $C_5 > 0$
such that for every $n$ with $h_n \leq h_\ast$ and every
$f \in L^2(V_n, m_n)$ with $\mathcal E_n(f) < \infty$,
\[
  \mathcal E_\infty(I_n f) \;\leq\; \mathcal E_n(f) +
    C_5\,h_n\,(\mathcal E_n(f) + \|f\|_{L^2(m_n)}^2).
\]
\end{lemma}

\begin{proof}[Sketch]
By Definition~\ref{def:transport}, $I_n f = \sum_v f(v) \psi_v$
with $\psi_v$ bounded partition of unity. Compute
$\nabla(I_n f) = \sum_v f(v) \nabla \psi_v$, which gives
\[
  |\nabla (I_n f)(p)|^2 \;=\;
    \biggl|\sum_v f(v) \nabla \psi_v(p)\biggr|^2
    \;\leq\; N_{\rm ov}\,\sum_v |f(v)|^2 |\nabla \psi_v(p)|^2.
\]
Since $\sum_v \psi_v \equiv 1$, $\sum_v \nabla \psi_v \equiv 0$.
Thus the constant mode of $f$ produces $\nabla I_n f = 0$, and
the gradient depends only on the differences $f(v) - f(u)$ via
$\nabla \psi_v$-localization to overlapping-support pairs. A
computation exactly parallel to Theorem~\ref{thm:form-consistency}
(Taylor expansion reversed, with the same isotropy constant
$8/3$ absorbed via $c_\Delta = 3/4$) gives
\[
  \int_{S^3} |\nabla I_n f|^2 d\mu_\infty \;=\; \mathcal E_n(f)
    + O(h_n)\cdot(\mathcal E_n(f) + \|f\|^2_{L^2(m_n)}).
\]
The proof mirrors Steps 1--4 of
Theorem~\ref{thm:form-consistency} with the roles of $R_n$
and $I_n$ exchanged; we omit the details.
\end{proof}

\begin{remark}[No BIK dependency]
\label{rmk:no-bik}
The proof is via Courant--Fischer min-max + Dirichlet-form
consistency (Theorem~\ref{thm:form-consistency}). We do
\emph{not} invoke the Burago--Ivanov--Kurylev theorem
\cite{BuragoIvanovKurylev}: their $\rho$-proximity-graph model
with $\varepsilon/\rho \to 0$ does not match our fixed
Schl\"{a}fli graph. The BIK paper is referenced only for
context. The Schl\"{a}fli-specific spectral-convergence theorem
presented here relies on the $W(D_4)$-induced isotropy
(Proposition~\ref{prop:v0-at-level-n},
Hypothesis~\ref{hyp:isotropy}), which is structurally different
from the BIK hypotheses.
\end{remark}

%=====================================================================
\section{Self-Audit}
\label{sec:audit}
%=====================================================================

This section separates unconditional from conditional content.

\textbf{Proved, conditional on Hypotheses~\ref{fact:distinct},
\ref{fact:tiling}, \ref{fact:separation}} (each with its own
status remark):
\begin{itemize}
\item[[1]] Riemannian-only target.
\item[[2]] $D_4$-branch only.
\item[[3]] $\rho_n = 1$, $c = 1$ printed.
\item[[4]] Limit $(S^3, g_{\rm round})$ printed exactly.
\item[[5]] Two metrics $d_n^{\rm ext}, d_n^{\rm int}$ defined
  explicitly (Def~\ref{def:dn}).
\item[[6]] $\iota_n$, $\mathcal X_n^{\rm sph}$ printed
  (Def~\ref{def:iota-n}, Hypothesis~\ref{fact:tiling}).
\item[[7]] \emph{Under Hypothesis~\ref{fact:tiling}:}
  $\varepsilon_n = O(h_n)$ net radius, $\iota_n$ isometric
  embedding of $(V_n, d_n^{\rm ext})$ into $(S^3, \dsph)$
  (Thm~\ref{thm:epsilon-net}).
\item[[8]] $\varepsilon$-net proved before GH.
\item[[9]] \emph{Under Hypothesis~\ref{fact:tiling}:} GH
  quantitative (Thm~\ref{thm:gh-convergence}).
\item[[10]] Measures defined (Def~\ref{def:measures}).
\item[[11]] $\Delta_n$ printed (Def~\ref{def:Delta-n}).
\item[[12]] $R_n, I_n$ bounded
  (Prop~\ref{prop:transport-bounded}).
\item[[15]] \emph{Under Hypothesis~\ref{fact:separation}:}
  doubling constants explicit (Prop~\ref{prop:doubling}).
\item[[17]] No cascade-gr.md, foundations.tex, scripts cited.
\item[[18]] No P4 cite.
\item[[19]] No Lorentzian/Einstein/NS/hydro.
\item[[20]] $S^3$ spectrum normalized to $c = 1$.
\item[[22]] Notation fixed, no drift.
\item $V_0$-orbit isotropy proved exactly
  (Prop~\ref{prop:v0-at-level-n}).
\end{itemize}

\textbf{Conditional PROGRAM (stated; not Clay-grade proofs; supporting
lemmas sketched, requiring a follow-up paper for full proofs):}
\begin{itemize}
\item Intrinsic-vs-extrinsic bi-Lipschitz comparison
  (Thm~\ref{thm:intrinsic-extrinsic}): $d_n^{\rm int} \leq
  d_n^{\rm ext} + C_2 h_n$ conditional on
  Hypothesis~\ref{hyp:isotropy}. The constant $C_2$ appears
  only here, in the conditional program; it does NOT appear in
  the base extrinsic theorems.
\item[[13]] Spectral convergence mode stated explicitly
  (Thm~\ref{thm:spectral-convergence}): as a conditional
  statement under Hypothesis~\ref{hyp:isotropy}, with proof via
  Courant--Fischer min-max using
  Lemma~\ref{lem:dual-form-consistency} and the quadrature
  Lemma~\ref{lem:quadrature}, both of which are sketched.
\item[[14]] Hypotheses for form consistency all stated.
\item[[16]] Form-convergence route via Courant--Fischer min-max.
\item[[21]] Eigenvalue convergence with multiplicity (conditional
  on Hypothesis~\ref{hyp:isotropy}).
\item \textbf{Status:} The conditional items above are a
  \emph{program} toward spectral convergence; they are not
  Clay-grade theorems until the sketched supporting lemmas are
  fully proved in a follow-up.
\end{itemize}

\textbf{What this paper does NOT do (honest acknowledgments):}
\begin{itemize}
\item The $O(h_n^2)$ distortion rate of earlier drafts is
  NOT proved; honest rate is $O(h_n)$.
\item The Burago--Ivanov--Kurylev 2014 theorem is NOT cited
  as load-bearing; their framework doesn't fit our graph.
\item Hypothesis~\ref{hyp:isotropy} for midpoint and centroid
  orbits is NOT proved here; deferred to follow-up.
\item The Dirichlet-form-consistency proof
  (Theorem~\ref{thm:form-consistency}) uses named external
  inputs (Poincar\'e, Sobolev embedding, Riemann-sum quadrature)
  whose hypotheses are verified but whose fully-inline proofs
  are not reproduced.
\end{itemize}

%=====================================================================
\section{Citation Contract}
\label{sec:handoff}
%=====================================================================

\textbf{Citable (extrinsic metric layer), conditional on the
structural hypotheses noted:}
\begin{itemize}
\item Theorem~\ref{thm:epsilon-net} (conditional on
  Hypothesis~\ref{fact:tiling}): $\varepsilon_n = O(h_n)$ net
  radius, plus $\iota_n$ is an isometric embedding of
  $(V_n, d_n^{\rm ext})$ into $(S^3, \dsph)$.
\item Theorem~\ref{thm:gh-convergence} (conditional on
  Hypothesis~\ref{fact:tiling}): $\dGH((V_n, d_n^{\rm
  ext}), (S^3, \dsph)) = O(h_n)$, with $h_n \leq M_\ast \alpha^n
  h_0$ for any $\alpha \in (1/\sqrt{3}, 1)$ up to absolute
  constant $M_\ast = M_\ast(\alpha)$.
\item Proposition~\ref{prop:doubling} (conditional on
  Hypotheses~\ref{fact:tiling} and~\ref{fact:separation}):
  doubling/packing for extrinsic-metric balls
  $B_{d_n^{\rm ext}}$.
\item Proposition~\ref{prop:v0-at-level-n} (unconditional):
  $V_0$-orbit isotropy constant $8/3$.
\end{itemize}

\textbf{Conditional PROGRAM (NOT Clay-grade citable as theorems;
supporting lemmas sketched; depend on Hypothesis~\ref{hyp:isotropy}
in addition to the structural hypotheses above):}
\begin{itemize}
\item Theorem~\ref{thm:intrinsic-extrinsic}: intrinsic-vs-
  extrinsic bi-Lipschitz comparison $d_n^{\rm int} \leq
  d_n^{\rm ext} + C_2 h_n$. The constant $C_2$ appears only
  here; conditional pearl-string proof sketched.
\item Theorem~\ref{thm:form-consistency}: Dirichlet-form
  consistency at rate $O(h_n\,\|f\|_{C^2}^2)$ for $f \in C^2(S^3)$.
  Proof uses the quadrature Lemma~\ref{lem:quadrature}, which is
  stated but whose proof is sketched.
\item Theorem~\ref{thm:spectral-convergence}: eigenvalue
  convergence via Courant--Fischer min-max. Proof uses the dual
  form-consistency Lemma~\ref{lem:dual-form-consistency}, which
  is stated but whose proof is sketched.
\end{itemize}

\textbf{Downstream citation guidance:} the conditional program
(Thms~\ref{thm:intrinsic-extrinsic},
\ref{thm:form-consistency}, \ref{thm:spectral-convergence} and
their supporting lemmas) is not Clay-grade citable until the
sketched proofs are completed in a follow-up paper. Cite the
extrinsic-metric items above as theorem-grade \emph{under} the
stated structural hypotheses (all of which are partially
verified: see Remarks~\ref{rmk:distinct-status},
\ref{rmk:tiling-status}, \ref{rmk:separation-status}).

%=====================================================================
\section{Notation Audit}
\label{sec:notation}
%=====================================================================

\begin{center}
\begin{tabular}{|l|l|l|}
\hline
Symbol & Meaning & First use \\
\hline
$V_n$ & level-$n$ $D_4$ vertex set &
  Notation~\ref{not:basic} \\
$x_v^{(n)}$ & P3/P5 Euclidean coordinate &
  Notation~\ref{not:basic} \\
$W(D_4)$ & $D_4$ Coxeter group & Notation~\ref{not:WD4} \\
$\iota_n$ & iterated spherical midpoint map &
  Definition~\ref{def:iota-n} \\
$h_n$ & mesh scale, $\leq M_\ast\,\alpha^n \pi/2$ for any $\alpha \in (1/\sqrt{3}, 1)$ &
  Definition~\ref{def:hn} \\
$e_n$, $c_e$ & minimum edge length and separation constant &
  Hypothesis~\ref{fact:separation} \\
$\ell_e^{(n)}$ & spherical edge length &
  Definition~\ref{def:dn} \\
$d_n^{\rm ext}$ & extrinsic metric $\dsph \circ \iota_n$ (base) &
  Definition~\ref{def:dn} \\
$d_n^{\rm int}$ & intrinsic graph metric (conditional program) &
  Definition~\ref{def:dn} \\
$\varepsilon_n$ & $O(h_n)$ net radius (under Hyp.~\ref{fact:tiling}) &
  Theorem~\ref{thm:epsilon-net} \\
$C_1$ & $\varepsilon$-net constant (under Hyp.~\ref{fact:tiling}) &
  Theorem~\ref{thm:epsilon-net} \\
$C_2$ & intrinsic-extrinsic distortion constant (conditional) &
  Theorem~\ref{thm:intrinsic-extrinsic} \\
$h_\ast$ & validity threshold & Theorem~\ref{thm:epsilon-net} \\
$m_n, \mu_\infty$ & normalized discrete and continuum measures &
  Definition~\ref{def:measures} \\
$c_0, C_0, C_D, C_P$ & doubling/packing constants &
  Proposition~\ref{prop:doubling} \\
$\pi_v$ & tangential projection in $T_{\iota_n(v)} S^3$ &
  Notation~\ref{not:pi-v} \\
$T_v^{(n)}$ & vertex-figure tensor &
  Definition~\ref{def:T-v} \\
$8/3$ & $V_0$-orbit isotropy constant &
  Prop~\ref{prop:v0-at-level-n} \\
$E_v^{(n)}$ & weighted-isotropy error tensor at $v$ &
  Hypothesis~\ref{hyp:isotropy} \\
$T_v^{(n),\,\rm w}$ & weighted vertex-figure tensor &
  Definition~\ref{def:T-v} \\
$\Delta_n$ & graph Laplacian &
  Definition~\ref{def:Delta-n} \\
$c_\Delta = 3/4$ & normalization &
  Definition~\ref{def:Delta-n} \\
$\mathcal E_n, \mathcal E_\infty$ & Dirichlet forms &
  Definition~\ref{def:Delta-n} \\
$\Delta_\infty = -\Delta_{S^3}$ & Laplace--Beltrami & \\
$R_n, I_n$ & transport operators &
  Definition~\ref{def:transport} \\
$\lambda_k^{(n)}, \lambda_k^\infty$ & $k$th eigenvalues &
  Theorem~\ref{thm:spectral-convergence} \\
$\dGH$ & Gromov--Hausdorff distance & abstract \\
\hline
\end{tabular}
\end{center}

\begin{thebibliography}{99}

\bibitem{AlgebraicSubstrate}
P2: \emph{Cascade Algebraic Substrate},
\texttt{papers/cascade-algebraic-substrate/cascade-algebraic-substrate.tex}.
Pins used: \S 3, \S 5.

\bibitem{RefinementSpaces}
P3: \emph{Refinement Spaces},
\texttt{papers/cascade-refinement-spaces/cascade-refinement-spaces.tex}.
Pins: Def.~2.1, Def.~2.2, Rmk.~2.3, Prop.~2.6.

\bibitem{MetricProjection}
P5: \emph{D_4 Metric Projection},
\texttt{papers/cascade-metric-projection/cascade-metric-projection.tex}.
Pins: Notation~2.1, Lem.~2.4.

\bibitem{SchlafliConvergence}
(This paper.)

\bibitem{BuragoBuragoIvanov}
D.~Burago, Yu.~Burago, S.~Ivanov, \emph{A Course in Metric
Geometry}, AMS GSM 33, 2001. Cor.~7.3.28.

\bibitem{BuragoIvanovKurylev}
D.~Burago, S.~Ivanov, Y.~Kurylev, \emph{A graph discretization of
the Laplace--Beltrami operator}, J.~Spectral Theory 4 (2014),
675--714. Cited for context only.

\bibitem{doCarmo}
M.~P.~do Carmo, \emph{Riemannian Geometry}, Birkh\"{a}user, 1992.

\bibitem{PetersenRiemannian}
P.~Petersen, \emph{Riemannian Geometry}, Springer, 2016.

\bibitem{BerardSpectralGeometry}
P.~H.~B\'{e}rard, \emph{Spectral Geometry}, LNM 1207, Springer,
1986.

\bibitem{CoxeterRegularPolytopes}
H.~S.~M.~Coxeter, \emph{Regular Polytopes}, Dover, 1973.

\bibitem{Humphreys}
J.~E.~Humphreys, \emph{Reflection Groups and Coxeter Groups},
Cambridge, 1990.

\bibitem{HeinonenAnalysis}
J.~Heinonen, \emph{Lectures on Analysis on Metric Spaces},
Springer, 2001.

\bibitem{ReedSimonI}
M.~Reed, B.~Simon, \emph{Methods of Modern Mathematical Physics
I}, Academic Press, 1980.

\bibitem{Kato}
T.~Kato, \emph{Perturbation Theory for Linear Operators},
Springer, 1995.

\bibitem{GilbargTrudinger}
D.~Gilbarg, N.~Trudinger, \emph{Elliptic Partial Differential
Equations of Second Order}, Springer, 1983.

\bibitem{BrennerScottFEM}
S.~C.~Brenner, L.~R.~Scott, \emph{The Mathematical Theory of
Finite Element Methods}, Springer, 2008.

\end{thebibliography}

\end{document}
